\section{Descent of the three local conditions for the large-prime count}

% BEGIN L79C1
\paragraph{L79C1. An integral retraction of the coefficient extension.}
At a place over $\ell$, write $K$ for the local base field, $L$ for
the chosen coefficient field and $L'/L$ for a finite coefficient
enlargement in the local comparison. Write $O,O'$ for their integer
rings. There is a continuous $L$-linear retraction $p:L'\to L$ with
$p(O')\subseteq O$. Indeed $O'$ is a finite torsion-free module over
the DVR $O$ and hence free. Its vector $1$ is primitive: if
$1\in\varpi O'$ for a uniformizer $\varpi$ of $O$, the element $1/\varpi$ would belong to $O'$ despite
its negative valuation. Thus the coordinates of $1$ in any free basis generate
the unit ideal. Choose an $O$-linear functional on that basis taking
$1$ to $1$ and extend by inverting $\varpi$. Finite-dimensionality
over $\mathbb Q_\ell$ gives continuity. Applying $p$ coordinatewise
commutes with every matrix of the representation over $L$ and
induces a retraction on $(L'/O')^a\to(L/O)^a$ for every rank $a$.
This is a map on the original lattice quotients, including when
$L'/L$ is ramified.

Let $z$ be a torsion cocycle for the original representation. If the
included cocycle is a coboundary $(g-1)v'$ on the local Galois
group, applying the quotient retraction gives $z(g)=(g-1)p(v')$.
The same argument, restricted to inertia, gives the native
unramified condition when the enlarged cocycle is a coboundary
there. In particular for a trivial ratio character it vanishes on
inertia. The local base field and inertia subgroup are held fixed
throughout this coefficient change.
% END L79C1

% BEGIN L79C2
\paragraph{L79C2. Descent of the Bloch--Kato finite condition.}
For the rank block suppose the lift of $z$ to the enlarged torsion
quotient belongs to the finite part $H^1_f$ in its cocycle formulation.
By the definition of the finite part on a quotient there are a
continuous rational cocycle $c'$ and a vector $m'\in (L')^r$ such that
the lift of $z(g)$ is represented by
$c'(g)+(g-1)m'$. Moreover $c'$ is crystalline as a cocycle:
$c'(g)\otimes1=(g-1)b'$ for one vector $b'$ in its tensor with
$B_{\rm cris}$, using the sign of the coboundary in that definition.
Indeed the quotient finite part is the sum of the reductions of
crystalline cocycles and torsion coboundaries. The first form a
submodule, and a primitive for a torsion coboundary lifts to a
rational vector.

Apply the coordinate retraction of L79C1. The projected cocycle
$c=p c'$ is continuous and the action commutes with $p$. The vector
$(p\otimes1)b'$ in the tensor over $\mathbb Q_\ell$ with
$B_{\rm cris}$ satisfies $(g-1)(p\otimes1)b'=c(g)\otimes1$.
On the quotient, the equality for $z(g)$ descends because integral
errors map to integral errors. Thus $c$ and $p(m')$ are precisely
the witnesses for $z\in H^1_f$ on the original lattice.
Combined with L79C1 this proves all three coefficient-descent
statements for the local cut cocycles: the rank part is finite
in the Bloch--Kato sense, the higher character ratios are
coboundaries, and the first character ratio is unramified, whenever
these properties hold after the given enlargement.

In particular if the enlarged conditions hold for every native cut
generator at $w$, their original $O$-span in the rank block lies in
the finite submodule, so the quotient by its intersection with that
submodule is zero. Its ordinary and normalized lengths are both zero.
The two character conclusions are pointwise on the full indicated
sets of cut cocycles. By G79E13 localization of each global generator
lies in those very cut spans. These deductions supply the local
conditions in the global count from the enlarged ones. The remaining
local input is their arithmetic construction there, using the exact
shifted local representations, their lattices and the local
Bloch--Kato criterion in the pinned proof.
% END L79C2
